\newif\ifglobal

%%%%%%%%%%%%%%%%%%%%%%%
%%% Compile to view %%%
%%%%%%%%%%%%%%%%%%%%%%%

%%%%%%%%%%%%%%%%%%%%%%%%%%%%%%%%%%%%%%%%%%%%%%%%%%%
%%% Readme:                                     %%%
%%% https://www.overleaf.com/read/bwdjvfjksvjy/ %%%
%%%%%%%%%%%%%%%%%%%%%%%%%%%%%%%%%%%%%%%%%%%%%%%%%%%

% >>> M?: marks a module setting number or important notice

% >>> M4: Set {./relative-path} to external files for this module <<<
\def \modulefiles {./13-PID}
% To add an external file, use:
% (Replace '---' with actual filename)
% '\input{\modulefiles/tables/---.tex}' for tables
% '\includegraphics{\modulefiles/figures/---}' for figures
% '\subfile{\modulefiles/---}' for register files

% >>> M1: This module is in global project? <<<
% 'YES' - uncomment; 'NO' - comment out
\globaltrue

\ifglobal
    \documentclass[main\_\_CN.tex]{subfiles}

\else
    % Font "FandolSong-Regular" used by ctex does not have GSUB table.
    % It causes warnings that do not affect compilation and can be ignored.
    \PassOptionsToPackage{quiet}{fontspec}

    \documentclass[openany, 10pt]{book}

    % >>> M2: In ./00-shared/preamble-trm-module, also find
    %         settings for visibility of todo notes, watermark <<<
    \usepackage{preamble-trm-module}


    % Enable Chinese typesetting
    \usepackage{ctex}

    % >>> M3: Temporary labels in Chinese
    %         you can add your own labels there <<<
    \myexternaldocument{./00-shared/temp-labels-cn}

\fi %\ifglobal


\setcounter{secnumdepth}{4}

% Define formatting of chapters in text
\setchineseformatting

% >>> M5: If this module needs any updates in
% './00-shared/preamble-trm-module.sty'
% (include additional package, etc.), contact your module PM
% or create the file './00-shared/changelog.tex' make the
% necessary updates yourself and document them in this file <<<

% Make text left-aligned and allow latex to break pages naturally, comment out for full justification
\raggedright
\raggedbottom


\begin{document}

\hypertarget{pidctrl}{}
\chapter{PID 控制器 (PID)}
\label{mod:pid}

\section{概述}
ESP32 双核芯片可以同时处理多个线程。PID 控制器在进程切换的过程中，
辅助完成进程号 PID 的切换。此外，PID 控制器还可以通过记录 CPU 处理中断的状态来管理嵌套中断。因此，PID 控制器可以使得用户在应用中更有效地管理进程切换和嵌套中断。

\section{主要特性}
PID 控制器有以下特性：
\begin{itemize}
    \item 管理进程优先级
    \item 切换进程号PID
    \item 记录中断信息
    \item 管理中断嵌套
\end{itemize}

\section{功能描述}
CPU 具有 8 个进程，进程号 PID 分别为 0 \verb+~+ 7。
这 8 个进程中，
PID 为 0/1 的进程比 PID 为 2 \verb+~+ 7 的进程拥有更多的权限。

CPU 在两种情形下会进行进程切换。
\begin{itemize}
    \item 当中断发生，CPU 从中断向量入口地址取指的时候。
无论中断发生之前正在运行的是哪个进程，
系统此时都将中断取指视为 PID 为 0 的进程。
  \item 当前进程主动切换进程时。
能够主动进行进程切换的进程一定是 PID 为 0/1 的高权限进程。
\end{itemize}

\subsection{中断识别}

CPU 一共有
Level 1、Level 2、Level 3、Level 4、Level 5、Level 6（Debug）、NMI
七个优先级的中断。
每一级中断有一个中断向量入口地址。
PID 控制器识别到 CPU 从中断向量入口地址取指的时候
自动切换 PID 至 0。
如果 CPU 只是对中断向量入口地址做数据访问，
那么 PID 控制器不采取任何动作。

PID 控制器能够识别的中断优先级取决于寄存器 \hyperref[regdesc:PIDCTRLINTERRUPTENABLEREG]{PIDCTRL\_INTERRUPT\_ENABLE\_REG}。
寄存器 \hyperref[regdesc:PIDCTRLINTERRUPTENABLEREG]{PIDCTRL\_INTERRUPT\_ENABLE\_REG} 中的某位为 1 时，
当 CPU 从此位对应的中断向量入口地址取指时，PID 控制器将切换进程；若寄存器 \hyperref[regdesc:PIDCTRLINTERRUPTENABLEREG]{PIDCTRL\_INTERRUPT\_ENABLE\_REG} 中的某位为 0 时，则 PID 控制器不采取任何动作。
这七级中断各自的中断向量入口地址由寄存器 PIDCTRL\_INTERRUPT\_ADDR\_1\_REG \verb+~+ PIDCTRL\_INTERRUPT\_ADDR\_7\_REG 决定，详见表 \ref{table:Interrupt}。


\begin{table}[H]
    \centering
    \caption{中断向量入口地址}
    \label{table:Interrupt}
    \begin{tabular}{|*{3}{c|}}
        \hline
        \rowcolor{lightgray}
       中断优先级 &
        \makecell{中断识别使能位\\PIDCTRL\_INTERRUPT\_ENABLE\_REG 位} &
       中断向量入口地址 \\
        \hline
        Level 1 & 1 & \hyperref[regdesc:PIDCTRLINTERRUPTADDR1REG]{PIDCTRL\_INTERRUPT\_ADDR\_1\_REG} \\
        \hline
        Level 2 & 2 &\hyperref[regdesc:PIDCTRLINTERRUPTADDR2REG]{PIDCTRL\_INTERRUPT\_ADDR\_2\_REG}\\
        \hline
        Level 3 & 3 & \hyperref[regdesc:PIDCTRLINTERRUPTADDR3REG]{PIDCTRL\_INTERRUPT\_ADDR\_3\_REG} \\
        \hline
        Level 4 & 4 & \hyperref[regdesc:PIDCTRLINTERRUPTADDR4REG]{PIDCTRL\_INTERRUPT\_ADDR\_4\_REG}\\
        \hline
        Level 5 & 5 & \hyperref[regdesc:PIDCTRLINTERRUPTADDR5REG]{PIDCTRL\_INTERRUPT\_ADDR\_5\_REG} \\
        \hline
        Level 6 ( Debug ) & 6 & \hyperref[regdesc:PIDCTRLINTERRUPTADDR6REG]{PIDCTRL\_INTERRUPT\_ADDR\_6\_REG} \\
        \hline
        NMI & 7 & \hyperref[regdesc:PIDCTRLINTERRUPTADDR7REG]{PIDCTRL\_INTERRUPT\_ADDR\_7\_REG} \\
        \hline
    \end{tabular}
\end{table}


\subsection{信息记录}

当 PID 控制器识别到中断时，
除了自动切换 PID 为 0，
还会记录 3 条信息：
\begin{enumerate}
    \item 当前中断的优先级
    \item 系统的上一次中断状态
    \item CPU 运行的上一个进程
\end{enumerate}
PID 控制器会将当前发生的中断的优先级
记录到寄存器 \hyperref[regdesc:PIDCTRLLEVELREG]{PIDCTRL\_LEVEL\_REG}，详见表 \ref{table:PIDCTRL_LEVEL_REG}。


\begin{table}[H]
    \centering
    \caption{PIDCTRL\_LEVEL\_REG}
    \label{table:PIDCTRL_LEVEL_REG}
    \begin{tabular}{|p{6cm}|p{8cm}|}
        \hline
        \rowcolor{lightgray}
       寄存器值 & 系统当前中断状态 \\
        \hline
        0 & 不处于中断中 \\
        \hline
        1 & 处于 Level 1 中断 \\
        \hline
        2 & 处于 Level 2 中断 \\
        \hline
        3 & 处于 Level 3 中断 \\
        \hline
        4 & 处于 Level 4 中断 \\
        \hline
        5 & 处于 Level 5 中断 \\
        \hline
        6 & 处于 Level 6 中断 \\
        \hline
        7 & 处于 NMI 中断\\
        \hline
    \end{tabular}
\end{table}


PID 控制器还将当前中断发生之前的状态
记录进寄存器 \hyperref[regdesc:PIDCTRLFROMnREG]{PIDCTRL\_FROM\_\textcolor{red}{\it{n}}\_REG}。
寄存器 \hyperref[regdesc:PIDCTRLFROMnREG]{PIDCTRL\_FROM\_\textcolor{red}{\it{n}}\_REG} 的位宽为 7。
其高 4 位表示此寄存器对应中断发生之前
系统的中断状态。
低 3 位表示此寄存器对应中断发生之前
系统处于哪个进程。
详见表 \ref{table:PIDCTRL_FROM_n_REG}。

\begin{table}[H]
    \centering
    \caption{PIDCTRL\_FROM\_\textcolor{red}{\it{n}}\_REG}
    \label{table:PIDCTRL_FROM_n_REG}
    \begin{tabular}{|p{1.5cm}|p{5.5cm}|}
        \hline
        \rowcolor{lightgray}
        \verb+[+6:3\verb+]+ & 当前中断发生前系统的中断状态 \\
        \hline
        0 & 不处于中断 \\
        \hline
        1 & 处于 Level 1 中断 \\
        \hline
        2 & 处于 Level 2 中断 \\
        \hline
        3 & 处于 Level 3 中断 \\
        \hline
        4 & 处于 Level 4 中断 \\
        \hline
        5 & 处于 Level 5 中断 \\
        \hline
        6 & 处于 Level 6 中断 \\
        \hline
        7 & 处于 Level 7 中断 \\
        \hline
    \end{tabular}
    \begin{tabular}{|p{1.5cm}|p{5.5cm}|}
        \hline
        \rowcolor{lightgray}
        \verb+[+2:0\verb+]+ & 当前中断发生前系统运行的进程 \\
        \hline
        0 & PID 为 0 的进程 \\
        \hline
        1 & PID 为 1 的进程\\
        \hline
        2 & PID 为 2 的进程\\
        \hline
        3 & PID 为 3 的进程\\
        \hline
        4 & PID 为 4 的进程\\
        \hline
        5 & PID 为 5 的进程\\
        \hline
        6 & PID 为 6 的进程\\
        \hline
        7 & PID 为 7 的进程\\
        \hline
    \end{tabular}
\end{table}


PID 控制器拥有寄存器 PIDCTRL\_FROM\_1\_REG \verb+~+ PIDCTRL\_FROM\_7\_REG，
它们分别对应
Level 1、Level 2、Level 3、Level 4、Level 5、Level 6（Debug）、NMI
七级中断。
系统通过这些寄存器处理中断嵌套，如图 \ref{figure:Interrupt_Nesting} 所示。

\begin{figure}[H]
    \centering
    \fbox{\includegraphics[width=0.9\textwidth]{\modulefiles/figures/Interrupt-Nesting_cn.png}}
    \caption{中断嵌套}
    \label{figure:Interrupt_Nesting}
\end{figure}


如果中断发生，
但是由于寄存器 \hyperref[regdesc:PIDCTRLINTERRUPTENABLEREG]{PIDCTRL\_INTERRUPT\_ENABLE\_REG} 的配置，
而使得 PID 控制器没有识别到它，
那么 PID 控制器不会做任何记录，
也不会更改寄存器 \hyperref[regdesc:PIDCTRLLEVELREG]{PIDCTRL\_LEVEL\_REG}、\hyperref[regdesc:PIDCTRLFROMnREG]{PIDCTRL\_FROM\_\textcolor{red}{\it{n}}\_REG} 的值。

\subsection{进程主动切换进程}

只有 PID 为 0/1 的进程可以主动切换进程。
切换之后的进程可以是 PID 为 0 \verb+~+ 7 中的任何一个。
进程主动切换进程的关键在于，
从当前进程的最后一条指令跳转到新进程的第一条指令的时候，
PID 要刚好从 0/1 变为新进程的 PID。

进程主动切换进程的软件流程为

\newcounter{STEP_MASK_NON_NMI}
\newcounter{STEP_MASK_NMI}
\newcounter{STEP_SET_DELAY}
\newcounter{STEP_SET_PID_NEW}
\newcounter{STEP_RESTORE}
\newcounter{STEP_CHANGE}
\newcounter{STEP_RELEASE_NON_NMI}
\newcounter{STEP_JUMP}
\begin{enumerate}
    \item 软件屏蔽除 NMI 中断之外的所有中断
        \setcounter{STEP_MASK_NON_NMI}{\theenumi}
    \item 将寄存器 \hyperref[regdesc:PIDCTRLNMIMASKENABLEREG]{PIDCTRL\_NMI\_MASK\_ENABLE\_REG} 置 1，生成 CPU NMI 中断屏蔽信号
        \setcounter{STEP_MASK_NMI}{\theenumi}
    \item 配置寄存器 \hyperref[regdesc:PIDCTRLPIDDELAYREG]{PIDCTRL\_PID\_DELAY\_REG}、\hyperref[regdesc:PIDCTRLNMIDELAYREG]{PIDCTRL\_NMI\_DELAY\_REG}
        \setcounter{STEP_SET_DELAY}{\theenumi}
    \item 配置寄存器 \hyperref[regdesc:PIDCTRLPIDNEWREG]{PIDCTRL\_PID\_NEW\_REG}
        \setcounter{STEP_SET_PID_NEW}{\theenumi}
    \item 配置寄存器 \hyperref[regdesc:PIDCTRLLEVELREG]{PIDCTRL\_LEVEL\_REG}、\hyperref[regdesc:PIDCTRLFROMnREG]{PIDCTRL\_FROM\_\textcolor{red}{\it{n}}\_REG}
        \setcounter{STEP_RESTORE}{\theenumi}
    \item 将寄存器 \hyperref[regdesc:PIDCTRLPIDCONFIRMREG]{PIDCTRL\_PID\_CONFIRM\_REG}、\hyperref[regdesc:PIDCTRLNMIMASKDISABLEREG]{PIDCTRL\_NMI\_MASK\_DISABLE\_REG} 置 1
        \setcounter{STEP_CHANGE}{\theenumi}
    \item 撤销对 NMI 中断之外的中断的屏蔽
        \setcounter{STEP_RELEASE_NON_NMI}{\theenumi}
    \item 切换到新进程取指
        \setcounter{STEP_JUMP}{\theenumi}
\end{enumerate}

虽然系统可以处理中断嵌套的情况，
但 PID 为 0/1 的进程不应在进程切换时被新的中断打断，
因此在步骤 \arabic{STEP_MASK_NON_NMI}、步骤 \arabic{STEP_MASK_NMI} 中屏蔽了所有中断。

步骤 \arabic{STEP_SET_DELAY} 中配置的
寄存器 \hyperref[regdesc:PIDCTRLPIDDELAYREG]{PIDCTRL\_PID\_DELAY\_REG}、\hyperref[regdesc:PIDCTRLNMIDELAYREG]{PIDCTRL\_NMI\_DELAY\_REG} 的值
将会作用于步骤 \arabic{STEP_CHANGE}。

步骤 \arabic{STEP_SET_PID_NEW} 中配置的
寄存器 \hyperref[regdesc:PIDCTRLPIDNEWREG]{PIDCTRL\_PID\_NEW\_REG} 的值
将会在步骤 \arabic{STEP_CHANGE} 之后成为新的进程 PID。

如果当前处于嵌套中断中且要恢复到上一个中断，
那么在步骤 \arabic{STEP_RESTORE} 中需要根据
寄存器 \textcolor{red}{\it{n}} 中记录的信息恢复寄存器 \hyperref[regdesc:PIDCTRLLEVELREG]{PIDCTRL\_LEVEL\_REG}。

步骤 \arabic{STEP_CHANGE} 中，将
寄存器 \hyperref[regdesc:PIDCTRLPIDCONFIRMREG]{PIDCTRL\_PID\_CONFIRM\_REG}、\hyperref[regdesc:PIDCTRLNMIMASKDISABLEREG]{PIDCTRL\_NMI\_MASK\_DISABLE\_REG} 置 1 后，
PID 控制器并不会立即将 PID 切换为 \hyperref[regdesc:PIDCTRLPIDNEWREG]{PIDCTRL\_PID\_NEW\_REG} 中的值，
也不会立即关闭 CPU NMI 中断屏蔽信号，
而是会分别等待一定数量的时钟周期数后才会执行这两个任务。
这两段等待时钟周期数即寄存器 \hyperref[regdesc:PIDCTRLPIDDELAYREG]{PIDCTRL\_PID\_DELAY\_REG}、\hyperref[regdesc:PIDCTRLNMIDELAYREG]{PIDCTRL\_NMI\_DELAY\_REG} 中的值。

在步骤 \arabic{STEP_RELEASE_NON_NMI} 中还可以执行其他任务，
只需要在步骤 \arabic{STEP_SET_DELAY} 中配置
寄存器 \hyperref[regdesc:PIDCTRLPIDDELAYREG]{PIDCTRL\_PID\_DELAY\_REG}、\hyperref[regdesc:PIDCTRLNMIDELAYREG]{PIDCTRL\_NMI\_DELAY\_REG} 时包含这些任务的时间开销即可。

\hypertarget{pidctrl-reg-summ}{}
\section{寄存器列表}\label{sec:pidctrl-reg-summ}

\subfile{\modulefiles/pidc_reg-summary_cn}

\section{寄存器}

本小节括号中的地址均为相对于 PID Controller 基地址的地址偏移量（相对地址），具体基地址见章节 \ref{mod:sysmem} \textit{\nameref{mod:sysmem}} 中的表 \ref{table:Peripheral_Address_Mapping} \textit{外设地址映射}。寄存器绝对地址见章节 \ref{sec:pidctrl-reg-summ} \textit{\nameref{sec:pidctrl-reg-summ}}。

\subfile{\modulefiles/pidc_reg_cn}

\end{document}
